\documentclass[11pt]{article}
\usepackage[margin=1in]{geometry}
\usepackage{amsmath}
\usepackage{booktabs}
\usepackage{url}
\setlength{\emergencystretch}{2em}
\title{A CLF-tube terminal set and recursive feasibility of the collision-avoidance controller}
\date{Prepared October 6, 2026}
\begin{document}
\maketitle

\section{Request}
The user asked for a terminal set that follows from the controller itself.
After the horizon, the controller is the same optimization without its PCBF
(collision) rows, so the vehicle dissipates to the nominal behavior under the
CLF. The terminal safe set is the set of states from which that controller
does not collide with a target obeying the motion contract (constant
tangential acceleration $A$ and sideslip $\beta$). The requirements were:
\begin{itemize}
\item no safety distance of its own;
\item no predesigned maneuver or mode;
\item continuous-time theory;
\item a target that leaves the 50-m perception range no longer exists;
\item ``no solution'' when the set cannot be reached.
\end{itemize}
Recursive feasibility was required without compromise. The user also asked
what the terminal set should be when the estimator and the controller are
designed together. The theory is in \path{controller/TERMINAL_SAFE_SET.md};
this report records the implementation and the measurements.

\section{Design}
\paragraph{CLF tube.} A terminal controller is any input law with
$\dot V\le-(2/T)V$ for the CLF $V=e^\top Pe$, $T=4$~s. Every error component
then obeys
$|e_i(t)|\le a_i\sqrt{V_0}\,e^{-t/T}$, $a_i=\sqrt{(P^{-1})_{ii}}$. The path
station stays within $s_0+v_{\mathrm{path}}t\pm G(t)$, where $G$ integrates a
bound on the path-speed error. The ego rectangle thus lies in a box family in
path coordinates (the tube) that depends only on $V_0$ and $s_0$.

\paragraph{Terminal set.} A state is terminal if all of the following hold:
\begin{itemize}
\item $V\le\bar c$;
\item its tube stays on the road;
\item its tube does not meet the target's forecast rectangle before the
  target leaves the range, within a 60-s encounter window.
\end{itemize}
$\bar c$ is the smaller of the CLF's certified region (1) and the largest
level whose ellipsoid lies inside the state rows. At 8~m/s it is 0.403 on
the straight road and 0.907 on the curve, set by rear adhesion; at 15~m/s it
is 1. The tube of a successor under the terminal controller lies inside the
tube of its predecessor shifted by the elapsed time. The set is therefore
forward invariant.

\paragraph{Recursive feasibility.} The issued plan is the nonlinear rollout of
$\text{anchor}+\alpha\cdot\text{correction}$ for the largest
$\alpha\in\{1,\tfrac12,\dots,\tfrac1{32},0\}$ that meets every hard row, ends
in the terminal set, and does not increase the prefix collision deficit. The
anchor is the previous accepted plan shifted by one hold. If that plan is
shorter than the prefix or its endpoint has left the set, it is extended by
terminal-controller holds. On the declared sampled model, with exact
information and the forecast unchanged, this anchor reproduces the previous
plan and is acceptable ($\alpha=0$). Hence no sample after the first accepted
one in an encounter reports no solution. The proposition has four
hypotheses: H1 (declared plant, exact state), H2 (forecast unchanged), H3
(the terminal controller has an admissible input, checked online) and H4
(the encounter ends within 60~s).

\paragraph{What changed.}
\begin{itemize}
\item \path{controller/terminalSafeSet.m} (new) implements the tube, the
  encounter exit, the bisected level, the state-row level, the terminal
  controller hold and the nonlinear row checks.
\item \path{solvePredictiveControl.m}:
  \begin{itemize}
  \item the anchor is the shifted accepted plan, extended by
    terminal-controller holds where needed;
  \item the startup rollout stops at its first node in the set;
  \item the separating and road-recovery rows are replaced by one cone,
    $\|F(e(y)+J\,\delta x_N)\|\le\sqrt{c^*}$;
  \item a step-size rule on the nonlinear rollout decides the issued plan,
    with at most two re-linearizations;
  \item the trust coefficient comes from the full step's nonlinear
    remainder.
  \end{itemize}
\item \path{collisionAvoidanceController.m}: state version 77 and new
  metadata.
\item \path{collisionAvoidanceControllerConfig.m}: \texttt{cfg.terminal}
  replaces \texttt{terminalSeparatingMarginMetersPerSecond} and
  \texttt{roadRecoveryAccelerationFraction}.
\end{itemize}
This reverses one earlier decision: the issued plan is now evaluated on the
nonlinear model before it is issued. Recursive feasibility needs the issued
plan to be a plan of the nonlinear sampled model, since only then does the
next frame's shift reproduce it. The step-size rule is part of the
optimization method; it never issues an input that violates the current
problem.

\section{Exact information}
\paragraph{Setup.} The 14 collision-threat encounters (7 scenarios at 8 and
15~m/s, \texttt{collisionThreatScenario}) were run with exact ego state and
exact target state inside the 50-m range. The runs used
\texttt{runNonlinearPredictiveSafetyValidation}, at most 2000 holds, with a
1-s recovery dwell after the last collision time. Two plants were used:
\begin{itemize}
\item ode45 (relative tolerance $10^{-11}$), the campaign default; it differs
  from the controller's one-step RK4 model by the integration error;
\item the declared RK4 model itself, on which hypothesis H1 holds exactly.
\end{itemize}
The per-run table is \path{TERMINAL_SAFE_SET_RECURSIVE_FEASIBILITY_20261006/runs.csv}.

\begin{center}
\begin{tabular}{lrr}
\toprule
 & ode45 plant & RK4 plant\\
\midrule
encounters recovered & 14 of 14 & 14 of 14\\
collisions; frames without a solution & 0; 0 & 0; 0\\
holds & 5630 & 5552\\
shifted frames & 5614 & 5536\\
shifted plan in the terminal set (certified anchor) & 5614 & 5536\\
shifted plan a feasible candidate ($\alpha=0$ acceptable) & 5607 & 5536\\
fresh restarts & 2 & 2\\
accepted full steps ($\alpha=1$) & 5188 & 5106\\
issued shifted plan ($\alpha=0$) & 28 & 41\\
largest horizon (holds) & 280 & 280\\
smallest continuous-time clearance (ode45 replay) & 0.155~m & 0.155~m\\
\bottomrule
\end{tabular}
\end{center}

\paragraph{Reading.} On the RK4 plant every one of the 5536 shifted plans was
a feasible candidate, as the proposition states. On the ode45 plant 7 of
5614 were not, all in the 8-m/s braking lead (frames 55, 63, 84--86, 101,
102), and all because of a collision row after the prefix.
\begin{itemize}
\item The accepted plan had that row active at the 0.2-m clearance.
\item The integration difference between ode45 and one-step RK4 then moved
  the row below the $10^{-6}$-m tolerance.
\item In these frames the optimizer's own step was accepted (6 full steps,
  1 half step). The candidate is a guarantee, not the only plan available.
\end{itemize}
The two fresh restarts are the second encounters on the circular road: the
oncoming target on the same circle re-enters the range after one synodic
period, at 77.4~s (8~m/s) and 54.4~s (15~m/s). A target that has left the
range does not exist for the controller, so its return is new information
and a new encounter, outside the proposition.

The smallest continuous-time clearance, 0.155~m (15-m/s accelerating
head-on), is below the 0.2-m clearance of the sampled rows. The rows hold at
nodes and hold midpoints only; between them the replayed clearance stayed
positive.

Controller time per frame had a 95th percentile of at most 0.37~s (ode45)
and 0.43~s (RK4), and a maximum of 1.05 and 1.14~s. These were measured with
up to 13 runs in parallel and memory swapping, and are not timing
measurements; the sample time is 0.05~s, so the controller is not real-time
in MATLAB either way.

The terminal level $\bar c$ binds in the plans: the optimizer has no reason
to keep the endpoint's $V$ low, so it rises to $\bar c$. A full step then often
lands just above it because of linearization error, which is why 13--15\% of
the steps were halved.

\section{Estimator in the loop}
\paragraph{What was run.} The noisy-sensor campaign had to stop early. It
uses the synthetic 0.1-s sensor contract and the NRMM adapter,
\texttt{estimatorControllerIntegrationConfig}, seeds 20261003 and 1--5. 84
runs were planned; only 5 ran on the final code, all at 8~m/s with seed
20261003.

Running 13 MATLAB sessions in parallel exhausted the 29~GB of memory, and the
kernel killed seven of them. After that, MathWorks online licensing refused
nearly every new session (``Unable to communicate with required MathWorks
services (error 5001)''). Network access to MathWorks was available, and a
lone \texttt{matlab -batch} start failed as well. The remaining 79 runs are
not reported here. The comparison with the previous design (82 of 84
recovered at \texttt{d1c31d6}) is therefore open.

\begin{center}
\begin{tabular}{lll}
\toprule
encounter (8~m/s, seed 20261003) & outcome & end\\
\midrule
head-on & recovered & 175 holds\\
crossing & recovered & 177 holds\\
accelerating head-on & no solution & hold 12 (\texttt{nonlinearAcceptanceFailed})\\
curved head-on & no solution & hold 9 (\texttt{nonlinearAcceptanceFailed})\\
braking lead & no solution & first frame (\texttt{pcbfNoNumericalResult})\\
\bottomrule
\end{tabular}
\end{center}
No run collided. In the 367 shifted frames before the failures, every anchor
was certified and every candidate feasible.

\paragraph{Diagnosis} (\path{noisy-failure-diagnosis.txt}, replayed from the
saved failure frames).
\begin{itemize}
\item \emph{Accelerating and curved head-on.} The shifted plan's endpoint is
  still in the terminal set under the new forecast, but a collision row after
  the prefix is violated: the forecast moved the target into the plan's tail.
  \begin{itemize}
  \item The convex problem became feasible only at larger trust scales (0.18
    and 0.66 of the box).
  \item Over the 85--95-hold horizon, the nonlinear rollout of the full step
    departed from its affine prediction by 0.47~m and 7.2~m.
  \item No step was acceptable, and the re-linearization did not recover.
  \end{itemize}
\item \emph{Braking lead, first frame.} The potential-field rollout ran
  512 holds without entering the terminal set and ended at $V=6.4$. At that
  station even the $V=0$ tube meets the target, so the terminal level is
  $-\infty$ and the startup problem is infeasible at every trust scale. With
  the first estimate $A=-1.1$~m/s$^2$, the lead's forecast stops and reverses
  down the lane, and the rollout never passes it.
\end{itemize}
The cause in both cases is the forecast, not the ego estimate. In the first
second of tracking, the estimated $A$ and $\beta$ swing between their
contract limits: from $+1.1$ to $-1.09$~m/s$^2$ and from $+0.055$ to
$-0.055$ within 0.3~s. The terminal set uses the forecast up to the target's
exit, several seconds ahead. The previous design issued the affine solution
without checking it on the nonlinear model, and its terminal condition was
local (separating at the endpoint), so these jumps did not make it fail at
these frames.

\paragraph{Size of the information change} (\path{noisy-joint-statistics.txt}).
The change of the forecast target position at the previous plan's endpoint
time, from one frame to the next, was:
\begin{center}
\begin{tabular}{lrrr}
\toprule
endpoint ahead of the frame & median & 95\% & max\\
\midrule
0--1~s & 0.07~m & 0.18~m & 0.35~m\\
1--2~s & 0.40~m & 0.79~m & 0.97~m\\
2--4~s & 1.00~m & 2.52~m & 5.24~m\\
4--8~s & 4.77~m & 17.5~m & 39.1~m\\
\bottomrule
\end{tabular}
\end{center}
The ego estimate departed from the plan's prediction by 0.0019~m (median),
0.0043~m (95\%) and 0.0064~m (max) in the body-point pose metric. In CLF
units ($\eta=\|F J(\hat x-x)\|$, truth known), the ego estimation error had a
95th percentile of 0.042 and a maximum of 0.047 on the straight road.

\section{Joint estimator--controller terminal set}
Section~9 of \path{controller/TERMINAL_SAFE_SET.md} derives three
requirements; the measurements above size them.
\begin{enumerate}
\item \emph{Ego: an ISS tube.}
  \begin{itemize}
  \item The tube radius becomes
    $r(t)=e^{-t/T}r_0+r_\infty(1-e^{-t/T})$, with
    $r_\infty\le\sqrt2\|F b\,\mathrm{diag}(\bar u)\|\,\eta/(1-\sqrt\rho)$.
  \item The constants are 16.4 and 16.5 at 8~m/s (straight, curve) and 23.4
    and 24.2 at 15~m/s (\path{iss-constants.txt}).
  \item The tube fits inside $\bar c$ only for $\eta\le0.039$ at 8~m/s on the
    straight road. The measured maximum, 0.047, is the actual error, not a
    certified bound, and already exceeds it.
  \item The levers are a faster CLF time constant (the factor scales with
    $T/h$; at the certified 2.1~s it halves) and terminal gains that are low
    in the poorly estimated directions.
  \end{itemize}
\item \emph{Plan: feedback instead of open loop.} The 7.2-m linearization gap
  of an 85-hold open-loop step is the open-loop propagation that Section~9.2
  describes. A plan executed as $u=v+K(x-z)$ about its affine prediction:
  \begin{itemize}
  \item keeps the convex problem unchanged;
  \item reproduces the old plan exactly under exact information;
  \item contracts both innovations and linearization error in the rollout.
  \end{itemize}
\item \emph{Target: nested information sets.} The forecast change grows
  from decimeters at 1~s to tens of meters at 4--8~s. A terminal set robust
  to it needs information sets that only shrink, as in a set-membership
  estimate of the constant $A$ and $\beta$ intersected over time. The
  published intervals stay at the contract width, so a terminal set robust
  to them is empty for exits more than a few seconds ahead. A
  constant-parameter set from the radar positions narrows roughly as
  $1/W^2$ for $A$ over an observation window $W$.
\end{enumerate}
Under the joint design, the terminal set is therefore (4.1) with the ISS
radius and the target's information set in place of its estimate. It is
nonempty when:
\begin{itemize}
\item the ego error satisfies $r_\infty^2<\bar c$;
\item the target set at the exit time is narrower than the lateral room the
  tube leaves.
\end{itemize}
Its recursive feasibility then needs nested target sets and a feedback-form
plan. None of the three is implemented; the present terminal set is
certainty-equivalent in the target and inflates the tube only by the ego
enclosure's position and yaw.

\section{Tests}
\texttt{runtests('tests')}: 590 of 590 passed on the final code, including
the new \texttt{terminalSafeSetTest} (invariance under the terminal
controller over 60 holds, nested levels, the state-row level against 2000
boundary points, the encounter window, the grid) and the shifted-plan
candidate test in \texttt{twoStagePredictiveControlTest}.

\section{Raw data}
The raw outputs under \path{~/.cache/collisionAvoidance/terminal-20261006}
are deleted after this report (standing instruction). The folder
\path{TERMINAL_SAFE_SET_RECURSIVE_FEASIBILITY_20261006} keeps the per-run
table, the statistics, the diagnosis log and the scripts.


\end{document}
